\ticket{13}{Методы поиска решений для игр с полной информацией: игровые деревья, минимаксный алгоритм, альфа-бета-отсечение}

\begin{examplan}
    \item Определить игры с полной информацией и игровые деревья.
    \item Объяснить минимаксный алгоритм.
    \item Описать ограничение глубины и оценочную функцию.
    \item Объяснить альфа-бета-отсечение и его эффективность.
\end{examplan}

\subsection*{Короткая версия для письменного ответа}

В играх с полной информацией текущее состояние и допустимые ходы известны всем игрокам. Для двух игроков с нулевой суммой поиск хода часто сводится к анализу \textbf{игрового дерева}, где узлы --- позиции, ребра --- ходы, а листья --- выигрыши, проигрыши или ничьи.

\subsection*{Игровые деревья}

Компоненты игрового дерева:
\begin{itemize}
    \item корень --- текущая или начальная позиция;
    \item узлы --- состояния игры;
    \item ребра --- возможные ходы;
    \item листья --- терминальные состояния;
    \item уровни обычно чередуют ходы MAX и MIN.
\end{itemize}

MAX стремится максимизировать оценку позиции, MIN --- минимизировать ее.

\subsection*{Минимакс}

\textbf{Минимакс} --- рекурсивный алгоритм выбора оптимального хода при предположении, что оба игрока действуют рационально.

\[
V(s)=
\begin{cases}
U(s), & \text{если } s \text{ терминальное состояние},\\
\max\limits_{s'\in Succ(s)} V(s'), & \text{если ход MAX},\\
\min\limits_{s'\in Succ(s)} V(s'), & \text{если ход MIN}.
\end{cases}
\]

Алгоритм:
\begin{enumerate}
    \item построить дерево до терминальных узлов или заданной глубины;
    \item оценить листья;
    \item поднять оценки вверх: максимум на узлах MAX, минимум на узлах MIN;
    \item выбрать ход из корня с лучшей оценкой.
\end{enumerate}

Если дерево полностью конечно и противник оптимален, минимакс дает оптимальную стратегию. Сложность по времени $O(b^m)$, где $b$ --- фактор ветвления, $m$ --- глубина; память при DFS --- $O(bm)$.

\subsection*{Ограничение глубины}

Для сложных игр полное дерево недоступно. Поэтому поиск ограничивают глубиной $d$, а нетерминальные позиции на этой глубине оценивают \textbf{эвристической оценочной функцией}. Например, в шахматах простая оценка может быть разностью материальной ценности фигур.

Сила игрового агента сильно зависит от качества оценочной функции и порядка рассмотрения ходов.

\subsection*{Альфа-бета-отсечение}

\textbf{Альфа-бета-отсечение} --- оптимизация минимакса, которая не меняет результат, но не просматривает ветви, не способные повлиять на итоговое решение.

\begin{itemize}
    \item $\alpha$ --- лучшее значение, которое MAX уже может гарантировать.
    \item $\beta$ --- лучшее значение, которое MIN уже может гарантировать.
\end{itemize}

Отсечение происходит, когда дальнейший просмотр узла бесполезен:
\begin{itemize}
    \item в MIN-узле можно отсекать, если текущее значение $v\leq\alpha$;
    \item в MAX-узле можно отсекать, если $v\geq\beta$.
\end{itemize}

В лучшем случае при хорошем упорядочивании ходов сложность снижается примерно с $O(b^m)$ до $O(b^{m/2})$, что позволяет просматривать дерево примерно вдвое глубже за то же время.

\begin{important}
    \item Игровое дерево описывает все возможные продолжения игры.
    \item Минимакс предполагает оптимальную игру обоих игроков.
    \item MAX выбирает максимум, MIN выбирает минимум.
    \item При ограниченной глубине нужна эвристическая оценочная функция.
    \item Альфа-бета-отсечение ускоряет минимакс, не меняя выбранный ход.
\end{important}
